\chapter{Speicherung und zeitabhängige Verfahren}\label{chap:dynamics}
\section{Konsistente Masse und Wärmekapazität}
$\rho_m$ und $c$ sind konstante positive Materialdaten pro Zelle. Aus dem
kinetischen Integral und dem thermischen Speicherterm folgen
\begin{equation}\label{eq:storage}
\mathbf S^e=\int_K\mathbf N^T\mathbf N\,\mathrm dV,\quad
\mathbf M^e=\rho_m(\mathbf S^e\otimes\mathbf I_3),\quad
\mathbf C_T^e=\rho_m c\mathbf S^e.
\end{equation}
Elektrische DOFs erhalten keine Masse. Mit den Momenten aus
Gleichung~\ref{eq:simplex-moment} ist für Tet4
\begin{equation}\label{eq:tet4-storage}
\mathbf S^e=\frac V{20}
\begin{pmatrix}2&1&1&1\\1&2&1&1\\1&1&2&1\\1&1&1&2\end{pmatrix},
\qquad\mathbf S_{\rm lumped}^e=\frac V4\mathbf I_4.
\end{equation}
Die Zeilensummenvariante wird ausdrücklich mit \texttt{lumped=True} gewählt.
Sie erhält Gesamtmasse je Translationsrichtung und Gesamtkapazität, nicht
alle Momente der konsistenten Matrix. Im Einheits-Tet4 mit $\rho_m=6$ kg/m$^3$
ist Gesamtmasse 1 kg, ein Diagonaleintrag 0,1 kg und ein Nebendiagonaleintrag
0,05 kg. Bei $c=2$ J/(kg K) beträgt Gesamtkapazität 2 J/K; gleichmäßiges
Heizen um 10 K benötigt 20 J.
\index{Massenmatrix}\index{Wärmekapazität}\index{HRZ-Lumping}

Bei Tet10 ergibt $\int N_i=-V/20$ an Ecken. Zeilensummen-Lumping würde
negative Eckmassen erzeugen. Angeboten wird nur die konsistente Matrix oder
explizites HRZ:
\begin{equation}\label{eq:hrz}
S_{{\rm HRZ},ii}=V\frac{S_{ii}}{\sum_jS_{jj}}.
\end{equation}
Die positive Diagonale wird auf das Gesamtvolumen skaliert; Masse/Kapazität
folgen durch $\rho_m$ bzw. $\rho_mc$. Dies entspricht der dokumentierten
Diagonalskalierung \cite{oasys-hrz}. Es ist keine exakte Integration aller
quadratischen Bewegungen. Ein Tet4-Booleschwert wird nicht still zu HRZ
umgedeutet. Unter-/Überlauf nicht darstellbarer Speicherkoeffizienten ist ein
Eingabefehler. Piezo-Dynamik verwendet die konsistente mechanische Masse.

\section{Transiente Wärme und Backward Euler}
Mit temperaturunabhängigen $\rho_m,c,\mathbf k$ gilt
\begin{equation}\label{eq:heat-time}
\rho_mc\,\dot T-\nabla\cdot(\mathbf k\nabla T)=Q,
\qquad\mathbf C_T\dot{\mathbf T}+\mathbf K_n\mathbf T=\mathbf f_n,
\end{equation}
wobei $\mathbf K_n$ Leit- und Robinanteil enthält. Ein vollständig aufgelöstes
Anfangsfeld $T_0$ und zeitliche Dirichlet-Daten sind nötig. Mit
$\Delta t_n=t_n-t_{n-1}>0$ wird
\begin{equation}\label{eq:backward-euler}
(\mathbf C_T/\Delta t_n+\mathbf K_n)\mathbf T_n
=\mathbf f_n+\mathbf C_T\mathbf T_{n-1}/\Delta t_n.
\end{equation}
Freie Zeilen mit bewegter Temperaturvorgabe $\mathbf g_n$ tragen zusätzlich
$-\mathbf K_{fc,n}\mathbf g_n-
\mathbf C_{T,fc}(\mathbf g_n-\mathbf g_{n-1})/\Delta t_n$.
Der volle Speicherterm darf nicht weggelassen werden, wenn feste Randknoten
Kapazität besitzen. Anfangswerte müssen Randwerte bei $t_0$ auf $10^{-10}$ K
erfüllen; es gibt keine stille Anfangsprojektion.
\index{Backward Euler}

Bekannte Signalereignisse werden in das Zeitraster aufgenommen. Quellen,
Flüsse und Robin-Daten verwenden am Intervallende den \emph{linken Grenzwert};
vorgegebene Temperaturen sind am Endpunkt rechtsstetig. Daher heizt ein
Quellensprung bei $t=2$ nicht schon das Intervall davor.
Callables müssen ihre Ereigniszeiten ausdrücklich melden; ein vorheriger
Float-Zeitpunkt approximiert den linken Grenzwert. Ein Raster mit endlichen,
streng steigenden, darstellbaren Schrittweiten ist erforderlich. Die robuste
Rampeninterpolation erweitert nicht die Grenzen des eigentlichen Zeitintegrators.

Für einen ungetriebenen skalaren Modus $\dot\vartheta=-\lambda\vartheta$ gilt
$\vartheta_n=\vartheta_{n-1}/(1+\lambda\Delta t)$.
Der Faktor ist für $\lambda\ge0$ und jeden positiven Schritt höchstens eins:
Backward Euler ist für dieses lineare dissipative Problem stabil, aber nur
zeitlich erster Ordnung und numerisch dämpfend. Stabilität bedeutet keine
zeitliche Auflösung. Konsistente Kapazität garantiert bei räumlicher
Diskretisierung keine allgemeine Temperaturpositivität/Monotonie.
Ein idealer Temperatursprung kann Überschwinger verursachen.

\subsection{Speicher- und Leistungsbilanz}
Die diskrete Änderung des thermischen Energieinhalts relativ zu $T_0$ ist
$\Delta U_T=\mathbf1^T\mathbf C_T(\mathbf T_n-\mathbf T_0)$ in J.
Das ist kein quadratischer Normausdruck. Pro Schritt gilt
\begin{equation}\label{eq:transient-balance}
P_Q+P_{D,\rm in}-P_{q,\rm out}-P_{h,\rm out}
-\mathbf1^T\mathbf C_T(\mathbf T_n-\mathbf T_{n-1})/\Delta t_n=0.
\end{equation}
Die Randreservoir-Reaktion enthält die Trägheit des Temperaturfeldes.
Feldresiduen und Bilanz werden unabhängig geprüft; initiale
Wärme-Frame-Reaktionen sind im Vertrag null, nicht als ausgewertete Bilanz
bei $t_0$ zu interpretieren.

\subsection{Handrechnung: gleichmäßige Abkühlung}
Ein im Test spezifiziertes gleichmäßiges Feld hat Rate $\lambda=1$/s,
$T_\infty=300$ K und $T_0=310$ K. Nach einem Schritt $\Delta t=0{,}1$ s gilt
$T_1=300+100/11=309{,}090909\ldots$ K.
Analytisch wäre $300+10e^{-0.1}=309{,}048374\ldots$ K.
Die Differenz ist Diskretisierungsfehler, keine Materialunsicherheit.
Auf 10/20/40/80 Schritten bis 1 s zeigen die bestehenden Tests Fehler
$0{,}176638/0{,}090100/0{,}045512/0{,}022873$ K und Ordnung gegen eins.
Raumkonvergenz wird separat mit zeitlich linearer Referenz geprüft, für die
Backward Euler keinen Zeitfehler erzeugt.

\section{Piezo-Dynamik als differential-algebraisches System}
Nach Elektroden- und Randreduktion sei $\mathbf D$ positiv definit. Für freie
mechanische/elektrische Unbekannte lautet das homogene Grundsystem
\begin{equation}\label{eq:dae}
\mathbf M\ddot{\mathbf u}+\mathbf K_u\mathbf u+\mathbf L\phi=\mathbf f_u,
\qquad\mathbf L^T\mathbf u-\mathbf D\phi=\mathbf f_\phi.
\end{equation}
Die mechanische starke Form ist $\rho_m\ddot{\mathbf u}-\nabla\cdot\sigma=\mathbf b$;
das Gauß-Gesetz bleibt algebraisch. Es gibt keine singuläre elektrische
Massenmatrix im Eigenwertsolver. Durch Auflösen der zweiten Gleichung folgt
\begin{align}\label{eq:dynamic-schur}
\phi&=\mathbf D^{-1}(\mathbf L^T\mathbf u-\mathbf f_\phi),\\
\mathbf K_{\rm eff}&=\mathbf K_u+\mathbf L\mathbf D^{-1}\mathbf L^T,\qquad
\mathbf M\ddot{\mathbf u}+\mathbf K_{\rm eff}\mathbf u
=\mathbf f_u+\mathbf L\mathbf D^{-1}\mathbf f_\phi.\nonumber
\end{align}
Inverse Schreibweise ist mathematische Kurznotation; Faktorisierung/Lösen
bzw. die dokumentierte dichte Schur-Referenz implementieren die Operation.
Bei bewegten mechanischen Rändern kommt $-\mathbf M_{fc}\ddot{\mathbf g}_c$
zur reduzierten Last; statische bekannte Randterme werden ebenfalls abgezogen.
Nur $\mathbf u_0,\mathbf v_0$ sind unabhängige Anfangsdaten.
$\phi_0$ folgt algebraisch und $\mathbf a_0$ aus der mechanischen Gleichung.
Explizite Anfangsdaten müssen vorgegebene Verschiebung/Geschwindigkeit bis
$10^{-12}$ m bzw. m/s erfüllen.
\index{Differential-algebraisches System}

\section{Modalanalyse und Randwirkung}
Homogene Störungen führen zu
\begin{equation}\label{eq:modes}
\mathbf K_{\rm eff}\mathbf z_j=\omega_j^2\mathbf M\mathbf z_j,
\qquad\mathbf z_i^T\mathbf M\mathbf z_j=\delta_{ij},
\qquad f_j=\omega_j/(2\pi).
\end{equation}
Potentialformen sind $\mathbf D^{-1}\mathbf L^T\mathbf z_j$.
Der größte mechanische Eintrag wird positiv gewählt; bei mehrfachen
Eigenwerten ist die Basis nicht eindeutig. Massennormierte Formen sind keine
physikalischen Anregungsamplituden. Nichtnull modale Randwerte werden
abgelehnt; ein statischer Bias verändert ohne Nichtlinearitätsmodell die
Tangenten nicht. Offene Leiter und kurzgeschlossene Leiter haben verschiedene
reduzierte Freiheitsräume und deshalb verschiedene Eigenfrequenzen.

Im Einheits-Tet4 mit nur einem freien mechanischen und einem elektrischen
DOF gilt $M=\rho_m/60$, $K=C/6$, $L=e/6$, $D=\kappa/6$.
Bei $\rho_m=6$, $C=10$, $e=2$, $\kappa=3$ in SI ist
$M=0{,}1$ kg, $K_{\rm eff}=17/9$ N/m und
$\omega_{\rm offen}^2=170/9$ s$^{-2}$.
Kurzschluss gibt $\omega_{\rm kurz}^2=50/3$ s$^{-2}$.
Die algebraische Rückabbildung lautet $\phi=(2/3)u$ in V/m mal u.

Für einen länglichen nichtaffinen dynamischen d33-Modus darf man Kurzschluss
nicht einfach durch $C$ ersetzen: Er fixiert nur die Gesamtspannung.
Die bestehende analytische Stabreferenz hat offen
$\omega=(\pi/(2L))\sqrt{C_D/\rho_m}$ und kurzgeschlossen
$\lambda\cot\lambda=e^2/(\kappa C_D)$,
$\omega=\lambda\sqrt{C_D/\rho_m}/L$. Diese eigene Referenz unterscheidet
statische uniforme Dehnung von einem räumlichen Schwingungsmodus.

\section{Newmark mit mittlerer Beschleunigung}
Die festen Parameter sind $\beta=1/4$, $\gamma=1/2$; physikalische und
numerische Dämpfung werden nicht ergänzt. Die üblichen kinematischen
Beziehungen \cite{opensees-newmark} lauten
\begin{align}\label{eq:newmark}
\mathbf u_{n+1}&=\mathbf u_n+\Delta t\mathbf v_n+
\Delta t^2[(1/2-\beta)\mathbf a_n+\beta\mathbf a_{n+1}],\\
\mathbf v_{n+1}&=\mathbf v_n+\Delta t[(1-\gamma)\mathbf a_n+
\gamma\mathbf a_{n+1}].\nonumber
\end{align}
Mit dem Prädiktor
$\mathbf u_p=\mathbf u_n+\Delta t\mathbf v_n+\Delta t^2\mathbf a_n/4$
setzt man $\mathbf u_{n+1}=\mathbf u_p+\Delta t^2\mathbf a_{n+1}/4$ in die
reduzierte Gleichung ein:
\begin{equation}\label{eq:newmark-solve}
(\mathbf M+\Delta t^2\mathbf K_{\rm eff}/4)\mathbf a_{n+1}
=\mathbf f_{{\rm eff},n+1}-\mathbf K_{\rm eff}\mathbf u_p.
\end{equation}
Für glatte lineare Daten ist die Methode auf festem Raumnetz zeitlich zweiter
Ordnung und ohne Dämpfung stabil; grobe Schritte haben weiterhin Phasenfehler.
Für $\ddot u+\omega^2u=0$ ist das Verfahren äquivalent zur
Trapezregel im Zustand $y=(u,v)$ mit
$\mathbf J=\left[\begin{smallmatrix}0&1\\-\omega^2&0\end{smallmatrix}\right]$:
\begin{equation}\label{eq:newmark-stability}
y_{n+1}=(\mathbf I-\Delta t\mathbf J/2)^{-1}
(\mathbf I+\Delta t\mathbf J/2)y_n,\qquad
g_\pm=\frac{1\pm i\omega\Delta t/2}{1\mp i\omega\Delta t/2}.
\end{equation}
$|g_\pm|=1$ für jeden positiven Schritt: Stabilität benötigt im linearen,
positiv definiten reduzierten Modell keine CFL-Schranke. Die numerische
Kreisfrequenz $\tilde\omega=2\arctan(\omega\Delta t/2)/\Delta t$
weicht aber um $-\omega^3\Delta t^2/12+O(\Delta t^4)$ ab. Energieerhaltung
beweist daher keine korrekte Phase. Für mehrere massennormierte Moden gilt
die Rechnung getrennt für jede Mode; nichtlineare oder instabile Modelle
werden dadurch nicht freigegeben.

Bekannte Kraft-/Spannungssprünge teilen das Raster: bis zum Ereignis wird links
integriert, danach werden Potential und Beschleunigung rechtsstetig neu
bestimmt. Verschiebung und Geschwindigkeit bleiben kontinuierlich.
Verschiebungs-/Geschwindigkeitssprünge sind nicht unterstützt.
\index{Newmark}

Bei ungetriebenen homogenen Rändern ist die positive Gesamtenergie
\begin{equation}\label{eq:dynamic-energy}
\mathcal E=\tfrac12\mathbf v^T\mathbf M\mathbf v+
\tfrac12\mathbf u^T\mathbf K_u\mathbf u+
\tfrac12\phi^T\mathbf D\phi
=\tfrac12\mathbf v^T\mathbf M\mathbf v+
\tfrac12\mathbf u^T\mathbf K_{\rm eff}\mathbf u
\end{equation}
konstant. Die zweite Gleichheit setzt die homogene elektrische Last voraus.
Bei Anregung wird keine ungeprüfte Arbeitsbilanz behauptet.
Für $u_0=0{,}01$ m im obigen Ein-DOF-Modell ist
$\mathcal E_0=17/180000$ J.
Die analytische Referenz $u(t)=0{,}01\cos(\sqrt{170/9}\,t)$ prüft auf
40/80/160 Schritten Fehlerquotienten 3,8--4,2 und Energietreue bis
$3\cdot10^{-12}$ relativ.

\section{Harmonische Analyse und Resonanzgrenze}
Spitzenwert-Phasoren verwenden
$\mathbf u(t)=\operatorname{Re}(\hat{\mathbf u}e^{+i\omega t})$,
$\hat{\mathbf v}=i\omega\hat{\mathbf u}$,
$\hat{\mathbf a}=-\omega^2\hat{\mathbf u}$.
$\omega=2\pi f$; Eingaben sind Hz, keine RMS-Werte. Dann
\begin{equation}\label{eq:harmonic}
(\mathbf K_{\rm eff}-\omega^2\mathbf M)\hat{\mathbf u}
=\hat{\mathbf f}_u+\mathbf L\mathbf D^{-1}\hat{\mathbf f}_\phi.
\end{equation}
Bewegte Randphasoren ergänzen $+\omega^2\mathbf M_{fc}\hat{\mathbf g}_c$
zusätzlich zu den statischen Randtermen. Komplexe Anteile bleiben erhalten.
$f=0$ reproduziert das statische System; negative Frequenzen sind ungültig.

Es gibt keine Dämpfung. Die Implementierung skaliert die dynamische Steifigkeit
mit $\operatorname{diag}(\mathbf K_{\rm eff}+\omega^2\mathbf M)^{-1/2}$.
Dimensionslose Eigenwertbeträge $\le10^{-10}$ führen zu einem expliziten
Resonanz-/Konditionsfehler, auch bei Last null. Das ist keine physikalische
Resonanzbreite und liefert keine künstlich gekappte Amplitude.
Ein Frequenzsweep kennt nur berechnete Frequenzen.
\index{Harmonik}

Im skalaren Beispiel und bei $\omega=1$ rad/s, $\hat F=1$ N,
$\hat f_\phi=0$ ist $\hat u=90/161$ m und $\hat\phi=60/161$ V.
Das große Verhältnis dient der linearen algebraischen Handrechnung;
physikalisch kleine Auslenkungen erhält man durch entsprechende Lastskalierung.

\section{Größen-, Genauigkeits- und Exportgrenzen}
Der dichte dynamische Schur-Pfad begrenzt \emph{aktive Gesamt-DOFs} vor der
Allokation auf 2048. Newmark-Historien sind auf
$\texttt{ntime}\cdot4\cdot\texttt{nactive}\le2\,000\,000$ beschränkt.
PETSc macht diese dichten Schritte nicht unbegrenzt skalierbar.
Wärme-Historien besitzen einen eigenen Byte-Grenzwert für numerische Arrays,
keine Garantie zum Gesamtprozess-RSS. Ereignisraster, räumlicher Fehler,
Zeitfehler und Frequenzabtastung werden separat beurteilt.
Zeit-/Moden-/Frequenzachsen sind im Export explizit; komplexe Felder haben
Realteil, Imaginärteil, Betrag und gültigkeitsmarkierte Phase.

\subsection*{Implementierung und Nachweise.}
{\small\raggedright
\path{src/fem_toolbox/elements/storage.py},
\path{src/fem_toolbox/solvers/transient_heat.py},
\path{src/fem_toolbox/solvers/piezo_dynamics.py},
\path{src/fem_toolbox/solvers/piezo_transient.py},
\path{src/fem_toolbox/solvers/piezo_harmonic.py};
\path{tests/unit/test_storage.py},
\path{tests/unit/test_transient_heat.py},
\path{tests/unit/test_piezo_modal.py},
\path{tests/unit/test_piezo_newmark.py},
\path{tests/unit/test_piezo_harmonic.py}.
Tet10-Modal-/Newmark-/Harmonik- und Raum-/Zeitnachweise sind in
Kapitel~\ref{chap:evidence} vollständig zugeordnet.

\par}
